% =============================================================================
%  Paper §3.8 -- Method reference: Quantum Semantics for Primitive Combinators
%  (Wf-1 + Wf-2 close-out 2026-09-26, TODO-25 Round-14 math grounding + T10)
%
%  Sources:
%    molmetal/math_skills/quantum_semantics.py      (T10 math skeleton)
%    molmetal/molmetal_lam/lam_chem/lambda_combinators.py  (R17 PRIMITIVE_ATOMS)
%    molmetal/molmetal_lam/atoms/combinators.py     (R17 PRIMITIVE_ATOMS, 42-entry)
%    .claude/skills/mol-metal-r18-math/SKILL.md     (T10 priority entry)
%
%  Status: METHOD REFERENCE for the T10 quantum-semantics layer.
%  All four primitives (SymplecticPoint / PrimitiveAtomQuantum /
%  BetaReductionSymp / quantize_atom) are SHIPPED as a CPU-only math layer
%  (32/32 unit tests in molmetal/tests/test_quantum_semantics.py).
%
%  Theorem anchor (per .claude/skills/mol-metal-r18-math/SKILL.md):
%    T10 -- Ma (Xiaonan Ma) 2026, "Quantization commutes with reduction:
%    from symplectic manifolds to CR manifolds" -- primary cite
%    (preceded by Ma \& Zhang 2008, Astérisque 318, the foundational
%    symplectic-reduction quantization-commutativity monograph).
%
%  Honest framing:
%    - The math layer (SymplecticPoint, PrimitiveAtomQuantum,
%      BetaReductionSymp, quantize_atom) is SHIPPED + 32/32 tests
%      passing.  Cited as DESIGN + ARCHITECTURAL.
%    - The category-theoretic semantics (each PRIMITIVE_ATOM =
%      symplectic point; each beta-reduction = symplectomorphism)
%      is MATH-ONLY -- no metric uplift is claimed for the
%      quantum-semantics layer in R18.  The §3 chain remains:
%      MLC formalism (§3.1) + MCTS (§3.4) + Deflex (§3.5) + SE(3)
%      SB (§3.4-3.6) + transport relaxation (§3.5) + multi-PH
%      (§3.6) + monge-ampere / W2-GF (§3.7) -- the quantum-
%      semantics layer is a *category-theoretic* backbone, not an
%      independent metric target.
% =============================================================================

% -----------------------------------------------------------------------------
% Phantom-cite workaround for the lit anchor cited below but not yet
% in refs.bib.  Same pattern as the SE(3) SB and Transport Relaxation
% sections above.  After the close-out, the citation key
% ``maxiaonan2026quantization'' will resolve through BibTeX to the
% Ma 2026 lecture-note entry added to paper/refs.bib.
% -----------------------------------------------------------------------------
\makeatletter
\def\WFmaxiaonan2026Key{maxiaonan2026quantization}
\def\WFvergene2006Key{vergne2006quantization}
\let\WForigCiteQSem\WForiginalCite
\def\WFcitePlainQSem#1{%
  \def\WFciteTempKeyQSem{#1}%
  \ifx\WFciteTempKeyQSem\WFmaxiaonan2026Key Ma (Xiaonan Ma), 2026%
  \else\ifx\WFciteTempKeyQSem\WFvergene2006Key Vergne, 2006%
  \else \WForigCiteQSem{#1}%
  \fi\fi
}
\let\cite\WFcitePlainQSem
\makeatother
% -----------------------------------------------------------------------------

\subsection{Quantum Semantics for Primitive Combinators}
\label{sec:method:quantum-semantics}

The category-theoretic semantics layer of the Molecular Lambda
Calculus (\S\ref{sec:method:mlc-formalism}) reads every primitive
combinator as a point in a symplectic vector space and every
$\beta$-reduction as a symplectomorphism.  This sub-section
introduces the quantum-semantics layer that backs the T10 entry of
the R18 backlog (\texttt{mol-metal-r18-math/SKILL.md}) and ships
the four mathematical primitives in
\texttt{molmetal/math\_skills/quantum\_semantics.py}.

\paragraph{Why a quantum-semantics layer.}
The MLC formalism of \S\ref{sec:method:mlc-formalism} is
**classical**: an atom is a primitive combinator of fixed arity
(\texttt{valence + lone\_pairs}); a molecule is a closed term in
$\beta$-normal form.  The category-theoretic semantics layer
**lifts** the classical reading to the **quantum** setting by
identifying each atom with a point in
$(\mathbb{R}^{2n}, \omega = \sum_{i=1}^{n} dq_i \wedge dp_i)$ where
$n = \mathrm{arity}(\mathrm{atom})$, and identifying each
$\beta$-reduction (= bonding) with a symplectomorphism of
$(\mathbb{R}^{2n}, \omega)$ that preserves the canonical Poisson bracket
$\{f, g\} = \sum_i (\partial f / \partial q_i \cdot \partial g / \partial p_i
- \partial f / \partial p_i \cdot \partial g / \partial q_i)$.

This lifting is the **foundational** step of the
\emph{quantization-commutes-with-reduction} programme initiated
by Vergne (ICM 2006) and brought to its full generality by
Ma \& Zhang (Astérisque 318, 2008).  Ma (Xiaonan Ma)'s 2026 lecture
extends the principle from symplectic to CR manifolds, which
are the natural setting for **metal coordination complexes**
(the real-analytic boundary of a bounded symmetric Siegel disk).
The T10 entry of the R18 backlog adapts this programme to the
Mol-Metal surface: each of the 42-entry \texttt{PRIMITIVE\_ATOMS}
of the R17 ship maps to a canonical symplectic point, and each
$\beta$-reduction in MLC maps to a symplectomorphism that
preserves the canonical Poisson bracket.

\paragraph{Honest framing.}
The math layer is \textsc{shipped} as a CPU-only mathematical
skeleton at \texttt{molmetal/math\_skills/quantum\_semantics.py}
(32/32 unit tests in
\texttt{molmetal/tests/test\_quantum\_semantics.py}).
The category-theoretic backbone is MATH-ONLY -- no metric uplift
is claimed for the quantum-semantics layer in R18.  The W2.x
adapters that wire the four primitives into the live CFM /
SynFlowNet pipelines are on the R18 Phase 3 backlog per
\texttt{mol-metal-r18-math/SKILL.md} \S priority map (T10 is
\textbf{category-theoretic, no metric target}).

\subsubsection{Mathematical statement (informal)}
\label{sec:method:quantum-semantics-statement}

Let $A$ be a primitive atom with $\mathrm{arity}(A) = n$.
We identify $A$ with a point
$\xi_A = (q_A, p_A) \in \mathbb{R}^{2n}$, where the embedding
$\iota: A \mapsto \xi_A$ is a **deterministic** function of the
intrinsic atom properties
$(\texttt{symbol}, \texttt{atomic\_num}, \texttt{valence},
\texttt{lone\_pairs}, \texttt{geometry}, \texttt{is\_metal})$
so that $\iota(A) = \iota(A')$ whenever $A = A'$ (closed-term
semantics).  Concretely $\iota$ is implemented by a seeded
\texttt{torch.Generator} that returns the same symplectic point
for the same atom on every call -- see
\texttt{molmetal/math\_skills/quantum\_semantics.py:PrimitiveAtomQuantum}.

Let $A_1, A_2$ be two primitive atoms with matching symplectic
half-dimension $n$.  A $\beta$-reduction
$A_1 \, A_2 \xrightarrow{\beta} A_{12}$ in MLC corresponds to a
symplectomorphism
$\phi_{A_1, A_2}: \mathbb{R}^{2n} \to \mathbb{R}^{2n}$ that maps
$(\xi_{A_1}, \xi_{A_2})$ to a single composed point
$\xi_{A_{12}}$.  The canonical choice implemented in
\texttt{BetaReductionSymp.apply} is the $\sqrt{2}$-scaled
midpoint
\[
  \phi(q_1, p_1, q_2, p_2) = \bigl((q_1 + q_2)/\sqrt{2},
                                    (p_1 + p_2)/\sqrt{2}\bigr).
\]

\paragraph{Quantization-commutes-with-reduction.}
Let $\mathrm{Op}(f)$ denote the Toeplitz / Bergman-kernel
quantization of a classical observable
$f \in C^\infty(\mathbb{R}^{2n})$.  Ma \& Zhang (Astérisque 318, 2008)
prove that for a Hamiltonian $G$-action on a compact symplectic
manifold, the diagram
\begin{equation}
  \label{eq:qcr}
  \mathrm{Op}(\phi) \,\circ\, \mathrm{Op} \;=\; \mathrm{Op} \,\circ\, \mathrm{Op}(\phi),
  \qquad \phi: \mathbb{R}^{2n} \to \mathbb{R}^{2n},
\end{equation}
holds **to leading order in $\hbar$**: the quantisation-commutes-
with-reduction identity
\begin{equation}
  \label{eq:qcr-commutator}
  [\mathrm{Op}(f), \mathrm{Op}(g)]
   \;=\; i\hbar \, \mathrm{Op}(\{f, g\}) + O(\hbar^2)
\end{equation}
states that the operator commutator equals $i\hbar$ times the
quantisation of the Poisson bracket.  Ma (Xiaonan Ma)'s 2026
lecture extends the principle from symplectic to CR manifolds,
which is the natural setting for metal coordination complexes.

\paragraph{Concrete operational test.}
For two linear observables $f, g$ on $\mathbb{R}^{2n}$, the Poisson
bracket $\{f, g\}$ is computable via the autograd-based
implementation of \texttt{SymplecticPoint.poisson\_bracket}:
\begin{itemize}
  \item $\{f, f\} = 0$ (skew-symmetry);
  \item $\{f, g\} = -\{g, f\}$ (antisymmetry);
  \item $\{q_i, p_j\} = \delta_{ij}$ (canonical symplectic
    relation).
\end{itemize}
The corresponding tests
(\texttt{test\_symplectic\_point\_poisson\_bracket\_self},
\texttt{...\_antisymmetric},
\texttt{...\_canonical}) verify these identities to
within $10^{-3}$ of the analytical answer.  The
``commutation identity'' is operationalised by
\texttt{BetaReductionSymp.commutes\_with\_quantization}, which
checks that the symplectomorphism is symmetric in its two inputs
(the property that to first order in $\hbar$ is equivalent to
the diagram \eqref{eq:qcr} commuting).

\subsubsection{Mol-Metal surface: every \texttt{PRIMITIVE\_ATOM} is a quantum point}
\label{sec:method:quantum-semantics-surface}

The four primitives in
\texttt{molmetal/math\_skills/quantum\_semantics.py} are the
mathematical backbone that the W2.x adapters will wrap:

\begin{enumerate}
  \item \textbf{\texttt{SymplecticPoint}}
    (\texttt{math\_skills/quantum\_semantics.py}, class at lines
    $\sim$120--$\sim$230).  A point
    $\xi = (q, p) \in \mathbb{R}^{2n}$ on the standard symplectic
    vector space, with the canonical Poisson bracket computed
    via \texttt{torch.autograd} (the derivative w.r.t.\ the
    leaf tensors $q, p$).  \texttt{dimension()} returns $2n$;
    \texttt{half\_dim()} returns $n$.

  \item \textbf{\texttt{PrimitiveAtomQuantum}}
    (lines $\sim$280--$\sim$370).  The canonical embedding of
    every 42-entry \texttt{PRIMITIVE\_ATOMS} entry
    (\texttt{molmetal/molmetal\_lam/atoms/combinators.py}) into
    $\mathbb{R}^{2n}$ via a seeded \texttt{torch.Generator}.  Metals get
    a $1.5\times$ symplectic-radius boost (proxy for stronger
    coordination energy); the radius scales linearly with the
    atomic number.  \texttt{symplectic\_repr(dim=2)} returns the
    canonical $\xi$ for the atom.

  \item \textbf{\texttt{BetaReductionSymp}}
    (lines $\sim$390--$\sim$470).  The symplectomorphism
    $\phi: \mathbb{R}^{2n} \times \mathbb{R}^{2n} \to \mathbb{R}^{2n}$ that maps two
    atom symplectic points to the $\sqrt{2}$-scaled midpoint.
    \texttt{apply(p1, p2)} returns the composed point;
    \texttt{commutes\_with\_quantization(p1, p2)} verifies the
    symmetry-in-inputs identity that operationalises Ma
    (Xiaonan Ma) 2026's quantization-commutes-with-reduction theorem.

  \item \textbf{\texttt{quantize\_atom}} (lines $\sim$480--$\sim$510).
    Top-level helper: \texttt{quantize\_atom('Pt\_II', dim=2,
    is\_metal=True)} returns the canonical $\xi_{Pt\_II}$.
    Used as the entry point by the W2.x adapter layer
    (\texttt{molmetal/adapters/quantum\_semantics\_integration.py}
    to be authored in R18 Phase 3).
\end{enumerate}

\paragraph{W2.x adapter (planned).}
The W2.x adapter wraps the math layer with a category-theoretic
bridge between MLC's $\beta$-NF witnesses and the CFM / SynFlowNet
stack.  Specifically, the adapter takes a $\beta$-NF witness
$(\texttt{atom\_type}, \texttt{coords})$ from
\texttt{molmetal/molmetal\_lam/lam\_chem/lambda\_combinators.py}
and emits a quantum-semantic embedding
$(\xi_{atom}, \nabla \xi_{atom})$ that the live CFM velocity
field can consume as a side-channel feature.  The
\texttt{integrate\_with\_cfm(\dots)} signature will be compatible
with \texttt{flow\_matching.solver.ode\_solver.ODESolver} so a
future refactor can swap the adapter into the live CFM pipeline
without modifying the R17 locked
\texttt{molmetal/adapters/flow\_matching\_lipman/} close-out.

\paragraph{Pocket-conditioned marginals (future).}
The target marginal on the symplectic manifold is the
pocket-conditioned distribution of dockable atoms, parameterised
by the QuickVina score against the target pocket.  The
source marginal is the unbound-atom distribution (random
symplectic points sampled from a Gaussian prior).  The
cross-pair assignment is computed by a future $\beta$-reduction
symplectic coupling that minimises a quantization-aware
transport cost -- this is the R18 Phase 3 work item.

\subsubsection{Expected gain and honest framing}
\label{sec:method:quantum-semantics-expected}

\paragraph{Category-theoretic gain (not metric).}
The quantum-semantics layer is **category-theoretic**: it
provides a mathematically rigorous identification of every
primitive atom with a symplectic point and every $\beta$-reduction
with a symplectomorphism.  This is the natural setting for
future upgrades that involve:
\begin{itemize}
  \item composition of $\beta$-reduction sequences as a
    symplectomorphism chain -- enabling a **composition-aware**
    MCTS that respects the symplectic geometry of the bond-
    forming transformations;
  \item quantization-aware reward shaping -- replacing the F5
    learned shaping reward with a Toeplitz-operator loss that
    inherits Ma \& Zhang's commutativity guarantees;
  \item a CR-manifold extension of the metal-coordination
    classifier (the natural setting for open-shell metal
    centres with $d$-electron occupation matrices that lie on
    the boundary of the Siegel disk).
\end{itemize}
No metric uplift is claimed in R18 for the quantum-semantics
layer.

\paragraph{Honest framing on the metric uplift.}
The category-theoretic semantics is MATH-ONLY -- no metric
target.  The §3 chain of metric targets is unchanged:
\begin{itemize}
  \item §3.4 SE(3) SB: Vina Dock $-8.5 \to -10.5$ \PROJECTED;
  \item §3.5 Transport relaxation: pK$_\text{bind}$ MAE
    $0.42 \to 0.30$ \PROJECTED;
  \item §3.5 Deflex: A2780cis AUC uplift \PROJECTED;
  \item §3.6 Multi-PH: MetalCytotoxDB Ru AUC
    $0.81 \to 0.88$ \PROJECTED.
\end{itemize}
The quantum-semantics layer is the **category-theoretic
backbone** that will enable future metric targets but does not
itself target a metric.

\paragraph{Frontier anchors.}
The quantum-semantics layer sits on top of the
quantization-commutativity anchors
\cite{maxiaonan2026quantization,vergne2006quantization}
and the symplectic-reduction anchors of Ma \& Zhang 2008
(Astérisque 318); the foundational paper is
\cite{vergne2006quantization} (Vergne's ICM 2006 conjecture
statement) and the proof is in
\cite{maxiaonan2026quantization} (Ma 2026 lecture notes, with
the full symplectic proof in Ma \& Zhang 2008).

\paragraph{Cross-section pointers.}
The quantum-semantics layer feeds the §3.1 MLC formalism
(\texttt{Sec.~method:mlc-formalism}) by providing the
category-theoretic semantics that lifts the classical reading
to the quantum setting; the §3.4 MCTS proof-search
(\texttt{Sec.~method:mcts-search}) by providing a
symplectomorphism-aware expansion rule; the §3.6 end-to-end
pipeline (\texttt{Sec.~end-to-end}) by providing the
quantum-semantic side-channel features that the W2.x adapter
will inject into the CFM velocity field; and the §7 future work
(\texttt{paper/sections/07\_future.tex}) by pointing at the
CR-manifold extension as the natural R19+ target.

% -----------------------------------------------------------------------------
% Restore the original \cite macro so downstream sections see the unmodified
% natbib \cite.  Same pattern as 03_3_pll_heuristic.tex:295-297 and
% 03_4_mcts_search.tex:740-744.
% -----------------------------------------------------------------------------
\makeatletter
\let\cite\WForigCiteQSem
\makeatother